%% ch13.tex --- Clustering with k-Means
%% Numbers are computed by verify/ch13_kmeans.py.
%% B60 and the snapshot run are the same data and run the companion lecture
%% animates frame by frame.  The six-point hand example (K6) is A-F below.

\chapter{Clustering with \texorpdfstring{$k$}{k}-Means}
\label{ch:ch13}

\epigraph{Birds of a feather flock together.}{---\ proverb}

\begin{chapterlearning}
After working through this chapter you will be able to:
\begin{enumerate}
  \item Explain what clustering is and why it is useful without labels.
  \item State the $k$-means algorithm and run it by hand.
  \item Define the distortion objective and prove that neither step of $k$-means
        can increase it.
  \item Explain why $k$-means converges, why only to a local minimum, and how
        empty clusters arise and are handled.
  \item Use multiple random initialisations and $k$-means\texttt{++} to avoid
        poor local minima.
  \item Choose the number of clusters by the elbow method --- and recognise when
        there is no elbow --- and from the cost of the application downstream.
  \item Apply $k$-means to image compression and compute the compression it
        achieves.
\end{enumerate}
\end{chapterlearning}

Every chapter so far has learned from labelled examples $(\vx, y)$. This part of
the book drops the labels. The data are $\xii{1},\dots,\xii{m}$ and nothing else,
and the task is to find structure in them.

The simplest such structure is a grouping. Customers who buy alike, documents on
the same topic, pixels of similar colour, genes with similar expression: in each
case we want to discover groups that no one has defined in advance. This chapter
develops the algorithm that is still the first choice for that task, analyses
exactly what it optimises, and shows how to use it well.

\section{What Clustering Is}
\label{sec:clustering}

\begin{defbox}{Clustering}
\term{Clustering} partitions a set of unlabelled examples into groups ---
\term{clusters} --- so that examples in the same cluster are similar to one
another and different from examples in other clusters. ``Similar'' must be
defined; here it is Euclidean distance between feature vectors.
\end{defbox}

Uses include market segmentation, grouping news articles by topic, organising
large image collections, finding families of related genes, and compressing data
by replacing each example with a representative of its cluster
(\cref{sec:compression}). Clustering is also a common first step in exploring a
new dataset, before anyone knows what the right labels would be.

\section{The \texorpdfstring{$k$}{k}-Means Algorithm}
\label{sec:kmeans-alg}

\begin{algorithmbox}{$k$-means}
\begin{lstlisting}[style=pseudocode]
choose K; initialise centroids mu_1, ..., mu_K
repeat until no assignment changes:
    # assignment step
    for i = 1 to m:
        c(i) <- index k of the centroid nearest x(i)
    # update step
    for k = 1 to K:
        mu_k <- mean of the x(i) with c(i) = k
\end{lstlisting}
Here $c^{(i)} \in \{1,\dots,K\}$ is the cluster of example $i$, and the
centroid $\mu_k \in \Real{n}$ is a point in feature space --- not, in general,
one of the examples. ``Nearest'' uses squared Euclidean distance
$\|\xii{i} - \mu_k\|^2$, which selects the same centroid as the distance itself
and avoids a square root.
\end{algorithmbox}

\begin{workedexample}{Six points by hand}
Cluster six points into $K = 2$ groups,
\[
  A(1,1),\ B(1.5,2),\ C(5,7),\ D(3.5,5),\ E(4.5,5),\ F(3.5,4.5),
\]
starting --- badly, on purpose --- with both centroids in the same corner:
$\mu_1 = A$, $\mu_2 = B$.

\textbf{Iteration 1, assign.} Squared distances to $(\mu_1, \mu_2)$:
$A$: $(0, 1.25)$; $B$: $(1.25, 0)$; $C$: $(52, 37.25)$; $D$: $(22.25, 13)$;
$E$: $(28.25, 18)$; $F$: $(18.5, 10.25)$. Only $A$ is nearer $\mu_1$, so
$c = (1, 2, 2, 2, 2, 2)$.

\textbf{Iteration 1, update.} $\mu_1 = A = (1, 1)$; $\mu_2$ is the mean of the
other five, $(3.6,\ 4.7)$.

\textbf{Iteration 2, assign.} $B$ is now at squared distance $1.25$ from $\mu_1$
and $11.70$ from $\mu_2$, so it switches: $c = (1, 1, 2, 2, 2, 2)$.

\textbf{Iteration 2, update.} $\mu_1 = (1.25,\ 1.5)$, the mean of $A$ and $B$;
$\mu_2 = (4.125,\ 5.375)$, the mean of $C$, $D$, $E$, $F$.

\textbf{Iteration 3, assign.} Nothing changes: converged. The cost of
\cref{sec:distortion} fell at every half-step:
\[
  13.0833 \;\to\; 3.3333 \;\to\; 1.5917 \;\to\; 1.0000 \;\to\; 1.0000 .
\]
\end{workedexample}

\begin{figure}[htbp]
  \centering
  \begin{tikzpicture}
    \begin{axis}[at={(0.0cm,0)}, width=4.1cm, height=4.1cm,
      xmin=0, xmax=6, ymin=0, ymax=8, title={\small start}, xtick=\empty, ytick=\empty]
      \addplot[only marks, mark=*, mark size=2pt, MLBlue] coordinates {(1,1)};
      \addplot[only marks, mark=square*, mark size=2pt, MLRose] coordinates {(1.5,2)(5,7)(3.5,5)(4.5,5)(3.5,4.5)};
      \addplot[only marks, mark=diamond*, mark size=4pt, draw=black, fill=MLBlue] coordinates {(1,1)};
      \addplot[only marks, mark=diamond*, mark size=4pt, draw=black, fill=MLRose] coordinates {(1.5,2)};
    \end{axis}
    \begin{axis}[at={(3.7cm,0)}, width=4.1cm, height=4.1cm,
      xmin=0, xmax=6, ymin=0, ymax=8, title={\small after iteration 1}, xtick=\empty, ytick=\empty]
      \addplot[only marks, mark=*, mark size=2pt, MLBlue] coordinates {(1,1)};
      \addplot[only marks, mark=square*, mark size=2pt, MLRose] coordinates {(1.5,2)(5,7)(3.5,5)(4.5,5)(3.5,4.5)};
      \addplot[only marks, mark=diamond*, mark size=4pt, draw=black, fill=MLBlue] coordinates {(1,1)};
      \addplot[only marks, mark=diamond*, mark size=4pt, draw=black, fill=MLRose] coordinates {(3.6,4.7)};
    \end{axis}
    \begin{axis}[at={(7.4cm,0)}, width=4.1cm, height=4.1cm,
      xmin=0, xmax=6, ymin=0, ymax=8, title={\small converged}, xtick=\empty, ytick=\empty]
      \addplot[only marks, mark=*, mark size=2pt, MLBlue] coordinates {(1,1)(1.5,2)};
      \addplot[only marks, mark=square*, mark size=2pt, MLRose] coordinates {(5,7)(3.5,5)(4.5,5)(3.5,4.5)};
      \addplot[only marks, mark=diamond*, mark size=4pt, draw=black, fill=MLBlue] coordinates {(1.25,1.5)};
      \addplot[only marks, mark=diamond*, mark size=4pt, draw=black, fill=MLRose] coordinates {(4.125,5.375)};
    \end{axis}
  \end{tikzpicture}
  \caption{The six-point example. Circles are assigned to cluster $1$, squares to
  cluster $2$; diamonds are centroids. The bad start puts five points in cluster
  $2$; one update moves $\mu_2$ into the upper group, and $B$ then switches
  sides.}
  \label{fig:k6}
\end{figure}

On larger data the same two steps take more iterations. \Cref{fig:b60-strip}
shows four moments of a run on B60, a dataset of $60$ points in three groups,
started from three randomly chosen examples. Two of the starting centroids fall
in the same group, and one centroid is dragged step by step across the top of
the plane until each group has its own; the run converges after eight
assignment steps, with the cost falling from $10.990$ to $2.338$.

\begin{figure}[htbp]
  \centering
  \begin{tikzpicture}
    \begin{axis}[at={(0.0cm,0.0cm)}, width=5.3cm, height=4.7cm,
      xmin=-1, xmax=11, ymin=0, ymax=11.3, xtick=\empty, ytick=\empty, title={\small iteration 1, $J = 10.990$}]
      \addplot[only marks, mark=*, mark size=1.3pt, MLBlue] coordinates {(2.30,6.32)(3.64,2.70)(3.86,7.49)(1.95,6.40)(-0.27,7.27)(2.56,7.71)(2.00,6.07)(3.94,7.64)(1.55,5.48)(3.15,7.54)(1.73,6.40)(4.29,4.42)(1.31,7.81)};
      \addplot[only marks, mark=square*, mark size=1.3pt, MLRose] coordinates {(8.42,10.47)(3.82,8.17)(1.13,8.00)(1.18,9.18)(0.68,9.61)(2.61,8.14)(2.04,8.45)(3.12,8.66)(7.89,9.87)(6.42,8.55)(1.57,10.23)(0.18,8.59)};
      \addplot[only marks, mark=triangle*, mark size=1.6pt, MLTeal] coordinates {(8.37,8.81)(8.11,8.28)(8.92,7.80)(7.79,6.04)(6.12,6.02)(4.05,2.91)(9.65,7.69)(7.59,5.90)(4.62,1.98)(4.99,2.42)(9.65,7.77)(5.60,3.24)(6.16,3.12)(10.10,8.26)(7.67,6.70)(4.68,2.83)(6.68,3.24)(7.85,7.66)(6.89,2.00)(6.69,6.84)(5.53,1.87)(4.26,2.44)(8.30,7.42)(5.43,0.91)(7.49,3.46)(6.57,2.67)(7.90,9.31)(5.02,2.03)(8.12,9.58)(5.72,3.48)(7.58,7.02)(4.78,1.85)(5.46,1.48)(8.62,7.17)(4.67,3.93)};
      \addplot[only marks, mark=diamond*, mark size=3.5pt, draw=black, fill=MLBlue] coordinates {(1.55,5.48)};
      \addplot[only marks, mark=diamond*, mark size=3.5pt, draw=black, fill=MLRose] coordinates {(1.57,10.23)};
      \addplot[only marks, mark=diamond*, mark size=3.5pt, draw=black, fill=MLTeal] coordinates {(7.49,3.46)};
    \end{axis}
    \begin{axis}[at={(5.6cm,0.0cm)}, width=5.3cm, height=4.7cm,
      xmin=-1, xmax=11, ymin=0, ymax=11.3, xtick=\empty, ytick=\empty, title={\small iteration 3, $J = 7.716$}]
      \addplot[only marks, mark=*, mark size=1.3pt, MLBlue] coordinates {(2.30,6.32)(1.13,8.00)(1.18,9.18)(0.68,9.61)(1.95,6.40)(-0.27,7.27)(2.56,7.71)(2.00,6.07)(1.55,5.48)(1.73,6.40)(0.18,8.59)(1.31,7.81)};
      \addplot[only marks, mark=square*, mark size=1.3pt, MLRose] coordinates {(8.42,10.47)(3.82,8.17)(2.61,8.14)(2.04,8.45)(3.86,7.49)(3.12,8.66)(3.94,7.64)(3.15,7.54)(7.89,9.87)(7.90,9.31)(6.42,8.55)(8.12,9.58)(1.57,10.23)};
      \addplot[only marks, mark=triangle*, mark size=1.6pt, MLTeal] coordinates {(8.37,8.81)(8.11,8.28)(8.92,7.80)(7.79,6.04)(6.12,6.02)(4.05,2.91)(9.65,7.69)(7.59,5.90)(4.62,1.98)(3.64,2.70)(4.99,2.42)(9.65,7.77)(5.60,3.24)(6.16,3.12)(10.10,8.26)(7.67,6.70)(4.68,2.83)(6.68,3.24)(7.85,7.66)(6.89,2.00)(6.69,6.84)(5.53,1.87)(4.26,2.44)(8.30,7.42)(5.43,0.91)(7.49,3.46)(6.57,2.67)(4.29,4.42)(5.02,2.03)(5.72,3.48)(7.58,7.02)(4.78,1.85)(5.46,1.48)(8.62,7.17)(4.67,3.93)};
      \addplot[only marks, mark=diamond*, mark size=3.5pt, draw=black, fill=MLBlue] coordinates {(1.74,6.90)};
      \addplot[only marks, mark=diamond*, mark size=3.5pt, draw=black, fill=MLRose] coordinates {(3.52,8.85)};
      \addplot[only marks, mark=diamond*, mark size=3.5pt, draw=black, fill=MLTeal] coordinates {(6.64,4.90)};
    \end{axis}
    \begin{axis}[at={(0.0cm,-5.0cm)}, width=5.3cm, height=4.7cm,
      xmin=-1, xmax=11, ymin=0, ymax=11.3, xtick=\empty, ytick=\empty, title={\small iteration 5, $J = 4.575$}]
      \addplot[only marks, mark=*, mark size=1.3pt, MLBlue] coordinates {(2.30,6.32)(3.82,8.17)(1.13,8.00)(1.18,9.18)(0.68,9.61)(2.61,8.14)(2.04,8.45)(3.86,7.49)(1.95,6.40)(-0.27,7.27)(2.56,7.71)(3.12,8.66)(2.00,6.07)(3.94,7.64)(1.55,5.48)(3.15,7.54)(1.73,6.40)(1.57,10.23)(0.18,8.59)(1.31,7.81)};
      \addplot[only marks, mark=square*, mark size=1.3pt, MLRose] coordinates {(8.37,8.81)(8.42,10.47)(8.11,8.28)(8.92,7.80)(9.65,7.69)(9.65,7.77)(10.10,8.26)(7.67,6.70)(7.85,7.66)(6.69,6.84)(8.30,7.42)(7.89,9.87)(7.90,9.31)(6.42,8.55)(8.12,9.58)(7.58,7.02)(8.62,7.17)};
      \addplot[only marks, mark=triangle*, mark size=1.6pt, MLTeal] coordinates {(7.79,6.04)(6.12,6.02)(4.05,2.91)(7.59,5.90)(4.62,1.98)(3.64,2.70)(4.99,2.42)(5.60,3.24)(6.16,3.12)(4.68,2.83)(6.68,3.24)(6.89,2.00)(5.53,1.87)(4.26,2.44)(5.43,0.91)(7.49,3.46)(6.57,2.67)(4.29,4.42)(5.02,2.03)(5.72,3.48)(4.78,1.85)(5.46,1.48)(4.67,3.93)};
      \addplot[only marks, mark=diamond*, mark size=3.5pt, draw=black, fill=MLBlue] coordinates {(1.61,7.70)};
      \addplot[only marks, mark=diamond*, mark size=3.5pt, draw=black, fill=MLRose] coordinates {(6.48,8.71)};
      \addplot[only marks, mark=diamond*, mark size=3.5pt, draw=black, fill=MLTeal] coordinates {(6.41,4.30)};
    \end{axis}
    \begin{axis}[at={(5.6cm,-5.0cm)}, width=5.3cm, height=4.7cm,
      xmin=-1, xmax=11, ymin=0, ymax=11.3, xtick=\empty, ytick=\empty, title={\small iteration 8, $J = 2.338$}]
      \addplot[only marks, mark=*, mark size=1.3pt, MLBlue] coordinates {(2.30,6.32)(3.82,8.17)(1.13,8.00)(1.18,9.18)(0.68,9.61)(2.61,8.14)(2.04,8.45)(3.86,7.49)(1.95,6.40)(-0.27,7.27)(2.56,7.71)(3.12,8.66)(2.00,6.07)(3.94,7.64)(1.55,5.48)(3.15,7.54)(1.73,6.40)(1.57,10.23)(0.18,8.59)(1.31,7.81)};
      \addplot[only marks, mark=square*, mark size=1.3pt, MLRose] coordinates {(8.37,8.81)(8.42,10.47)(8.11,8.28)(8.92,7.80)(7.79,6.04)(6.12,6.02)(9.65,7.69)(7.59,5.90)(9.65,7.77)(10.10,8.26)(7.67,6.70)(7.85,7.66)(6.69,6.84)(8.30,7.42)(7.89,9.87)(7.90,9.31)(6.42,8.55)(8.12,9.58)(7.58,7.02)(8.62,7.17)};
      \addplot[only marks, mark=triangle*, mark size=1.6pt, MLTeal] coordinates {(4.05,2.91)(4.62,1.98)(3.64,2.70)(4.99,2.42)(5.60,3.24)(6.16,3.12)(4.68,2.83)(6.68,3.24)(6.89,2.00)(5.53,1.87)(4.26,2.44)(5.43,0.91)(7.49,3.46)(6.57,2.67)(4.29,4.42)(5.02,2.03)(5.72,3.48)(4.78,1.85)(5.46,1.48)(4.67,3.93)};
      \addplot[only marks, mark=diamond*, mark size=3.5pt, draw=black, fill=MLBlue] coordinates {(2.02,7.76)};
      \addplot[only marks, mark=diamond*, mark size=3.5pt, draw=black, fill=MLRose] coordinates {(8.09,7.86)};
      \addplot[only marks, mark=diamond*, mark size=3.5pt, draw=black, fill=MLTeal] coordinates {(5.33,2.65)};
    \end{axis}
  \end{tikzpicture}
  \caption{Four moments in a $k$-means run on B60 (the companion lecture shows
  every step as an animation). Shapes show the current assignment, diamonds the
  centroids. At the start, blue and rose share the upper-left group and teal
  covers the other two; by iteration $5$ the rose centroid has moved most of the
  way to the upper-right group; at iteration $8$ the clusters have $20$ points
  each and the assignment no longer changes.}
  \label{fig:b60-strip}
\end{figure}

\begin{pitfall}{Empty clusters}
If no example is nearest some centroid, that cluster's mean is $0/0$. Start all
three centroids of B60 bunched in one corner --- at $(1,1)$, $(2,1.5)$ and
$(1.5,3)$ --- and exactly this happens at the first update. The two standard
remedies are to drop the empty cluster and continue with $K-1$, or to
re-initialise its centroid at a random example (or at the example farthest from
its centroid). Library implementations do the latter automatically.
\end{pitfall}

\section{The Distortion Objective}
\label{sec:distortion}

\begin{defbox}{Distortion}
The \term{distortion} (cost) of an assignment and a set of centroids is
\[
  J\bigl(c^{(1)},\dots,c^{(m)},\mu_1,\dots,\mu_K\bigr)
  = \frac1m\sum_{i=1}^{m}\bigl\|\xii{i} - \mu_{c^{(i)}}\bigr\|^2 ,
\]
the average squared distance from each example to the centroid of its cluster.
\end{defbox}

$k$-means is not a heuristic that happens to produce clusters. It is an
optimizer for $J$.

\begin{derivation}{Neither step of $k$-means can increase $J$}
$k$-means is \term{coordinate descent}: it minimises $J$ over one block of
variables with the other held fixed, then alternates.

\textbf{Assignment step} (centroids fixed, choose $c$). $J$ is a sum of
independent terms $\|\xii{i} - \mu_{c^{(i)}}\|^2$, one per example, and each is
minimised by choosing the nearest centroid. So the assignment step finds the
best possible $c$ for the current centroids, and cannot increase $J$.

\textbf{Update step} (assignment fixed, choose $\mu$). The terms involving
$\mu_k$ are $\sum_{i:c^{(i)}=k}\|\xii{i} - \mu_k\|^2$, a convex quadratic in
$\mu_k$. Its gradient is
\[
  \frac{\partial}{\partial\mu_k}\sum_{i:c^{(i)}=k}\|\xii{i} - \mu_k\|^2
  = -2\sum_{i:c^{(i)}=k}\bigl(\xii{i} - \mu_k\bigr),
\]
which is zero exactly when $\mu_k = \frac{1}{|C_k|}\sum_{i\in C_k}\xii{i}$, the
mean of the cluster. So the update step finds the best possible centroids for
the current assignment, and cannot increase $J$ either.
\end{derivation}

\begin{keyinsight}{Convergence, and its limits}
$J$ never increases and is bounded below by $0$. There are finitely many
possible assignments --- at most $K^m$ --- and the algorithm cannot revisit one
without $J$ having stayed constant, so it terminates in finitely many steps. A
cost that rises from one step to the next is therefore proof of a bug.

What $k$-means cannot do is guarantee the \emph{global} minimum. Coordinate
descent stops at any configuration where neither step can improve $J$, and on
most datasets there are many such configurations.
\end{keyinsight}

\begin{researchspotlight}{Lloyd (1982) --- Least squares quantization in PCM}
\textbf{The problem.} In pulse-code modulation, a continuous signal is
represented by a finite set of levels. Which levels minimise the mean squared
error between the signal and its quantised version?

\textbf{The idea.} Alternate two conditions that any optimal quantiser must
satisfy: each value is mapped to its nearest level, and each level is the mean
(centroid) of the values mapped to it. Iterating them improves the quantiser
monotonically.

\textbf{The result.} A practical algorithm for optimal scalar quantisation,
written at Bell Labs in 1957 and circulated internally for twenty-five years
before formal publication. Applied to vectors of features rather than samples of
a signal, it is $k$-means, which is why the algorithm is often called
\emph{Lloyd's algorithm}.

\textbf{Why it is here.} It is \cref{sec:kmeans-alg}, and its framing as
quantisation is exactly the image compression of \cref{sec:compression}.
\parencite{lloyd1982}
\end{researchspotlight}

\section{Initialisation and Local Minima}
\label{sec:km-init}

\begin{defbox}{Random initialisation}
Choose $K < m$ training examples at random and use them as the initial
centroids.
\end{defbox}

On B60, $200$ runs from random initialisations end as follows:

\begin{center}
\small
\renewcommand{\arraystretch}{1.12}
\begin{tabular}{@{}lr@{}}
\toprule
final distortion $J$ & runs (of $200$)\\
\midrule
$2.338$ (the three true groups)  & $193$\\
$7.717$ to $8.152$ (local minima) & $7$\\
\bottomrule
\end{tabular}
\end{center}

In each poor run one centroid covers two of the groups while two centroids split
the third --- a configuration neither step can repair, because no single point is
nearer a different centroid.

\begin{algorithmbox}{Multiple random initialisations}
\begin{lstlisting}[style=pseudocode]
for t = 1 to T:              # typically 50 to 1000
    randomly initialise k-means
    run it to convergence
    record the clustering and its distortion J
return the clustering with the lowest J
\end{lstlisting}
If one run finds the best clustering with probability $p$, all $T$ runs miss it
with probability $(1-p)^T$. On B60, $p = 193/200 = 0.965$, and two runs miss with
probability $0.035^2 \approx 0.0012$.
\end{algorithmbox}

\subsection{\texorpdfstring{$k$}{k}-means++}

Random initialisation often places two starting centroids close together, as in
\cref{fig:b60-strip}. \term{$k$-means++} spreads them out deliberately: choose
the first centroid uniformly at random from the data; then choose each next
centroid at random with probability proportional to its squared distance from the
nearest centroid already chosen. A point far from every existing centroid is
much more likely to be picked. In one draw on B60, the farthest point from the
first centre has probability $0.0447$ of being chosen next, against $1/60 =
0.0167$ under uniform sampling.

Over $200$ runs on B60, $k$-means++ reaches the best clustering in $197$, against
$193$ for random initialisation. B60 is an easy problem; on data with many
clusters the difference is much larger, and $k$-means++ comes with a guarantee
random initialisation lacks.

\begin{researchspotlight}{Arthur \& Vassilvitskii (2007) --- $k$-means++}
\textbf{The problem.} $k$-means is fast but can be arbitrarily bad: from a poor
start, its final distortion can be unboundedly worse than the optimum.

\textbf{The idea.} Seed the centroids by $D^2$ sampling --- each new centroid
chosen with probability proportional to its squared distance from the nearest
existing one --- then run ordinary $k$-means.

\textbf{The result.} The expected distortion of the seeding alone is within a
factor of $O(\log K)$ of the optimum, before any iterations; in experiments the
method was both faster and more accurate than random seeding.

\textbf{Why it is here.} It is the initialisation of \cref{sec:km-init}, and the
default in \tool{scikit-learn}'s \code{KMeans}, which also runs ten
initialisations and keeps the best.
\parencite{arthur2007}
\end{researchspotlight}

\section{Choosing \texorpdfstring{$K$}{K}}
\label{sec:choose-k}

The distortion always decreases as $K$ grows, given good initialisations, and
reaches $0$ when $K = m$. So $K$ cannot be chosen by minimising $J$.

\begin{defbox}{The elbow method}
Plot the best distortion found against $K$, and look for an \term{elbow}: a value
of $K$ after which further clusters reduce $J$ much more slowly.
\end{defbox}

\begin{figure}[htbp]
  \centering
  \begin{tikzpicture}
    \begin{axis}[mlaxis, width=0.72\linewidth, height=5.0cm,
      xlabel={number of clusters $K$}, ylabel={distortion $J$},
      xmin=0.5, xmax=8.5, ymin=0, ymax=18.5, xtick={1,...,8},
      legend style={font=\scriptsize, at={(0.98,0.97)}, anchor=north east}]
      \addplot[MLBlue, very thick, mark=*, mark size=1.6pt] coordinates
        {(1,14.407)(2,8.132)(3,2.338)(4,1.923)(5,1.519)(6,1.269)(7,1.056)(8,0.805)};
      \addlegendentry{B60: three groups}
      \addplot[MLAmber, very thick, dashed, mark=square*, mark size=1.4pt] coordinates
        {(1,17.351)(2,9.128)(3,5.567)(4,3.962)(5,3.023)(6,2.435)(7,1.998)(8,1.646)};
      \addlegendentry{U60: uniform, no groups}
    \end{axis}
  \end{tikzpicture}
  \caption{Elbow curves, each point the best of $50$ random starts. B60 has a
  sharp elbow at $K = 3$, its true number of groups. Sixty points scattered
  uniformly over a square have no groups, and their curve bends smoothly with no
  elbow at all.}
  \label{fig:elbow}
\end{figure}

\begin{pitfall}{The elbow that isn't}
On B60 the elbow is unmistakable: $J$ falls from $14.41$ to $8.13$ to $2.34$,
then to only $1.92$. But the uniform data of \cref{fig:elbow} also produces a
decreasing curve, and someone determined to find an elbow will find one
somewhere in it. Real datasets usually look more like the uniform curve than
like B60. An elbow is evidence of structure when it is sharp; its absence is
common, and is not a failure of the method.
\end{pitfall}

\begin{keyinsight}{Choose $K$ for the purpose}
Clustering is usually a step towards something else, and that something
usually settles $K$. A clothing maker clustering customers' heights and weights
into sizes might choose $K = 3$ (S, M, L) or $K = 5$ (XS to XL): five sizes fit
customers better, three are cheaper to make and stock. The data cannot decide
that trade-off; the business can. Evaluate candidate values of $K$ by how well the
resulting clusters serve the downstream task.
\end{keyinsight}

\section{Application: Image Compression}
\label{sec:compression}

A colour image stores $24$ bits per pixel: $8$ each for red, green and blue.
Treat every pixel as a point in $\Real{3}$, cluster the pixels into $K$ colours,
and replace each pixel by its cluster's centroid. The image now needs only
$\lceil\log_2 K\rceil$ bits per pixel, plus a small table of the $K$ colours.

\begin{workedexample}{Compressing a photograph}
Take the $427\times640$ colour photograph distributed with \tool{scikit-learn}
(\code{load\_sample\_image("china.jpg")}): $273{,}280$ pixels, $800.6$~KiB at $24$
bits each. Run $k$-means on a random sample of $5000$ pixels with colour values
scaled to $[0,1]$, then assign every pixel to its nearest centroid:

\begin{center}
\small
\renewcommand{\arraystretch}{1.12}
\begin{tabular}{@{}rrrrr@{}}
\toprule
$K$ & bits per pixel & size (KiB) & compression & mean squared colour error\\
\midrule
$2$  & $1$ & $33.4$  & $24.0\times$ & $0.05929$\\
$4$  & $2$ & $66.7$  & $12.0\times$ & $0.02112$\\
$8$  & $3$ & $100.1$ & $8.0\times$  & $0.00976$\\
$16$ & $4$ & $133.5$ & $6.0\times$  & $0.00547$\\
$32$ & $5$ & $166.9$ & $4.8\times$  & $0.00324$\\
\bottomrule
\end{tabular}
\end{center}

Each doubling of $K$ costs one bit per pixel and roughly halves the colour error
--- the elbow trade-off of \cref{sec:choose-k}, now with a price in storage.
\end{workedexample}

\begin{figure}[htbp]
  \centering
  \begin{tikzpicture}
    \fill[fill={rgb,255:red,13;green,13;blue,7}, draw=MLGray!50] (0.00,0.00) rectangle ++(0.95,0.95);
    \node[font=\tiny, MLGray] at (0.475,-0.18) {8.3\%};
    \fill[fill={rgb,255:red,39;green,30;blue,22}, draw=MLGray!50] (1.15,0.00) rectangle ++(0.95,0.95);
    \node[font=\tiny, MLGray] at (1.625,-0.18) {8.3\%};
    \fill[fill={rgb,255:red,49;green,53;blue,45}, draw=MLGray!50] (2.30,0.00) rectangle ++(0.95,0.95);
    \node[font=\tiny, MLGray] at (2.775,-0.18) {6.0\%};
    \fill[fill={rgb,255:red,84;green,64;blue,31}, draw=MLGray!50] (3.45,0.00) rectangle ++(0.95,0.95);
    \node[font=\tiny, MLGray] at (3.925,-0.18) {5.1\%};
    \fill[fill={rgb,255:red,80;green,83;blue,69}, draw=MLGray!50] (4.60,0.00) rectangle ++(0.95,0.95);
    \node[font=\tiny, MLGray] at (5.075,-0.18) {5.3\%};
    \fill[fill={rgb,255:red,114;green,100;blue,54}, draw=MLGray!50] (5.75,0.00) rectangle ++(0.95,0.95);
    \node[font=\tiny, MLGray] at (6.225,-0.18) {5.0\%};
    \fill[fill={rgb,255:red,111;green,114;blue,101}, draw=MLGray!50] (6.90,0.00) rectangle ++(0.95,0.95);
    \node[font=\tiny, MLGray] at (7.375,-0.18) {4.6\%};
    \fill[fill={rgb,255:red,156;green,125;blue,79}, draw=MLGray!50] (8.05,0.00) rectangle ++(0.95,0.95);
    \node[font=\tiny, MLGray] at (8.525,-0.18) {3.2\%};
    \fill[fill={rgb,255:red,145;green,151;blue,142}, draw=MLGray!50] (0.00,-1.45) rectangle ++(0.95,0.95);
    \node[font=\tiny, MLGray] at (0.475,-1.63) {4.6\%};
    \fill[fill={rgb,255:red,218;green,153;blue,113}, draw=MLGray!50] (1.15,-1.45) rectangle ++(0.95,0.95);
    \node[font=\tiny, MLGray] at (1.625,-1.63) {1.9\%};
    \fill[fill={rgb,255:red,175;green,185;blue,182}, draw=MLGray!50] (2.30,-1.45) rectangle ++(0.95,0.95);
    \node[font=\tiny, MLGray] at (2.775,-1.63) {5.1\%};
    \fill[fill={rgb,255:red,202;green,208;blue,213}, draw=MLGray!50] (3.45,-1.45) rectangle ++(0.95,0.95);
    \node[font=\tiny, MLGray] at (3.925,-1.63) {7.5\%};
    \fill[fill={rgb,255:red,191;green,214;blue,240}, draw=MLGray!50] (4.60,-1.45) rectangle ++(0.95,0.95);
    \node[font=\tiny, MLGray] at (5.075,-1.63) {8.9\%};
    \fill[fill={rgb,255:red,230;green,231;blue,233}, draw=MLGray!50] (5.75,-1.45) rectangle ++(0.95,0.95);
    \node[font=\tiny, MLGray] at (6.225,-1.63) {5.1\%};
    \fill[fill={rgb,255:red,217;green,234;blue,252}, draw=MLGray!50] (6.90,-1.45) rectangle ++(0.95,0.95);
    \node[font=\tiny, MLGray] at (7.375,-1.63) {7.8\%};
    \fill[fill={rgb,255:red,242;green,246;blue,252}, draw=MLGray!50] (8.05,-1.45) rectangle ++(0.95,0.95);
    \node[font=\tiny, MLGray] at (8.525,-1.63) {13.4\%};
  \end{tikzpicture}
  \caption{The $16$ colours $k$-means found for the photograph, ordered from dark to
  light, with the share of pixels assigned to each. They are the centroids: every
  pixel of the compressed image is one of these colours. The palette is dominated
  by greys, browns and the pale blues of the sky.}
  \label{fig:palette}
\end{figure}

\begin{labbox}{Lab 13 --- $k$-means}
Implement \code{find\_closest\_centroids} and \code{compute\_centroids}, each
vectorised, and the loop that alternates them. Reproduce the six-point example of
\cref{sec:kmeans-alg}, including the distortion after every half-step, and
confirm that it never increases. Run $200$ random initialisations on B60 and
reproduce the table of \cref{sec:km-init}; then implement $k$-means++ and compare.

Finally compress an image: load it, cluster its pixels with $K = 16$, rebuild the
image from the palette, and report the bits per pixel, the compression factor and
the mean squared colour error.
\end{labbox}

\section*{Chapter Summary}
\addcontentsline{toc}{section}{Chapter Summary}

\begin{itemize}
  \item Clustering finds groups in unlabelled data; $k$-means does so by
        alternating two steps: assign each example to its nearest centroid, and
        move each centroid to the mean of its examples.
  \item $k$-means is coordinate descent on the distortion
        $J = \frac1m\sum_i\|\xii{i} - \mu_{c^{(i)}}\|^2$. The assignment step is
        optimal for fixed centroids, the update step optimal for a fixed
        assignment, so $J$ never increases and the algorithm terminates.
  \item It converges to a local minimum, not necessarily the global one: $7$ of
        $200$ random starts on B60 stop at $J \approx 8$ instead of $2.338$.
  \item Run many random initialisations and keep the lowest $J$; $k$-means++
        seeds centroids by $D^2$ sampling and does better still.
  \item Empty clusters arise from bad starts; drop or re-initialise them.
  \item $J$ always falls with $K$. Use a sharp elbow when there is one --- B60's
        is at $K = 3$ --- and the downstream purpose always.
  \item Clustering pixel colours compresses images: $16$ colours store a
        photograph in a sixth of the space with mean squared colour error
        $0.0055$.
\end{itemize}

\section*{Exercises}
\addcontentsline{toc}{section}{Exercises}

\exercisesection{Short Questions}

\sq Distinguish clustering from classification.

\sq State the two steps of $k$-means.

\sq Why is squared distance used in the assignment step?

\sq What is an empty cluster, and how is it handled?

\sq Define the distortion $J$.

\sq Why does $k$-means always terminate?

\sq Why can $k$-means converge to different clusterings from different starts?

\sq Describe $k$-means++ initialisation.

\sq Why can $K$ not be chosen by minimising $J$?

\sq How does $k$-means compress an image, and how many bits per pixel does
$K = 64$ need?

\exercisesection{Analytical Problems}

\ap Run $k$-means by hand on the six points of \cref{sec:kmeans-alg} starting from
$\mu_1 = A$, $\mu_2 = C$. Give every assignment, centroid and distortion.

\ap Prove that the mean minimises $\sum_{i\in C_k}\|\xii{i} - \mu\|^2$ over $\mu$,
both by setting the gradient to zero and by expanding the square around the mean.

\ap Prove that the assignment step cannot increase $J$, and that the update step
cannot either. Hence prove that $k$-means terminates.

\ap Construct a dataset of four points in the plane and an initialisation for
which $k$-means with $K = 2$ converges to a clustering that is not the global
minimum. Compute both distortions.

\ap If a single random start finds the best clustering with probability $0.3$,
how many starts are needed for a $99\%$ chance of finding it?

\ap Show that if the medians replace the means in the update step, the resulting
algorithm decreases $\frac1m\sum_i\|\xii{i} - \mu_{c^{(i)}}\|_1$ instead. Why is
this \emph{$k$-medians} algorithm more robust to outliers?

\exercisesection{Design Problems}

\dpr A supermarket chain wants to segment its loyalty-card customers from a year
of purchases. Design the features (and their scaling), the choice of $K$, and how
the marketing team will use and evaluate the segments.

\dpr $k$-means on $10^7$ points with $K = 100$ is too slow to restart $500$
times. Design a cheaper procedure that still guards against poor local minima,
and say how you would check it.

\dpr Design three small datasets with known answers that would expose,
respectively, a wrong assignment rule, a wrong centroid update, and missing
empty-cluster handling in a $k$-means implementation.

\dpr Design an experiment that tells a client whether the clusters you report are
real structure or an artefact of random initialisation.

\exercisesection{Programming Exercises}

\pe Implement vectorised $k$-means and reproduce the six-point example and the
$200$-run table for B60.

\pe Implement $k$-means++ and compare it with random initialisation on a dataset
of your choice with $K = 10$ well-separated groups.

\pe Reproduce the elbow curves of \cref{fig:elbow} for B60 and U60.

\pe Compress an image of your choice with $K = 2, 4, 8, 16, 32$ and plot the mean
squared colour error against bits per pixel.

\exercisesection{Self-Test}

\qz The update step of $k$-means sets each centroid to: (a)~the nearest example;
(b)~the mean of its assigned examples; (c)~the median; (d)~a random example.

\qz The distortion $J$ during $k$-means: (a)~can increase; (b)~never increases;
(c)~is always zero; (d)~oscillates.

\qz $k$-means is guaranteed to find: (a)~the global minimum; (b)~a local minimum;
(c)~the true number of clusters; (d)~labels.

\qz $k$-means++ chooses each new centroid with probability proportional to:
(a)~its distance; (b)~its squared distance to the nearest chosen centroid;
(c)~a constant; (d)~its label.

\qz As $K$ increases, the best distortion: (a)~increases; (b)~decreases;
(c)~stays the same; (d)~is undefined.

\qz Compressing an image to $16$ colours needs per pixel: (a)~$16$ bits;
(b)~$8$ bits; (c)~$4$ bits; (d)~$2$ bits.

\qz If a centroid receives no examples: (a)~$J$ becomes negative; (b)~its mean is
undefined and it must be dropped or re-initialised; (c)~the algorithm has
converged; (d)~$K$ doubles.

\furtherreading

\textcite{james2021} §12.4.1 is the primary reading: it states the $k$-means
objective as the within-cluster variation, proves that the algorithm decreases it
at every step, and illustrates local optima with six runs from different starts
--- the experiment of \cref{sec:km-init}. \textcite{bishop2006} §9.1 develops
$k$-means as coordinate descent on the distortion measure, applies it to image
segmentation and compression with figures well worth seeing, and leads into §9.2,
where $k$-means appears as a limiting case of fitting a mixture of Gaussians ---
the probabilistic view behind the anomaly detection of \cref{ch:ch14}.

\bigskip
\noindent\emph{Next:} \Cref{ch:ch14} turns from groups to outliers: modelling
what normal data looks like, and flagging what does not.
